\DocumentMetadata{
  pdfversion=2.0,
  pdfstandard=ua-2,
  lang=en-US,
  tagging=on,
  tagging-setup={math/setup=mathml-SE,tabsorder=structure}
}
\documentclass[11pt,letterpaper]{article}
\usepackage[margin=0.68in,includefoot]{geometry}
\usepackage{newcomputermodern}
\usepackage{amsmath,amscd,mathtools}
\usepackage{tikz,adjustbox}
\providecommand{\square}{\mdlgwhtsquare}
\usepackage[hidelinks]{hyperref}
\hypersetup{
  pdftitle={Topology Qualifying Exam, January 2026, Part 1},
  pdfauthor={University of Florida Department of Mathematics},
  pdfsubject={Graduate mathematics examination},
  pdfdisplaydoctitle=true,
  bookmarksopen=true
}
\setcounter{secnumdepth}{0}
\setlength{\parindent}{0pt}
\setlength{\parskip}{0.28em}
\setlength{\emergencystretch}{3em}
\setlength{\leftmargini}{2.15em}
\setlength{\leftmarginii}{2.65em}
\setlength{\leftmarginiii}{3em}
\renewcommand{\labelenumi}{\textbf{\theenumi.}}
\pagestyle{empty}
\raggedbottom
\allowdisplaybreaks
\newenvironment{examproblems}[1]{%
  \begin{enumerate}%
  \setcounter{enumi}{#1}%
  \setlength{\topsep}{0.35em}%
  \setlength{\partopsep}{0pt}%
  \setlength{\itemsep}{0.48em}%
  \setlength{\parsep}{0.18em}%
}{\end{enumerate}}
\newenvironment{examsubparts}{%
  \begin{enumerate}%
  \setlength{\topsep}{0.18em}%
  \setlength{\partopsep}{0pt}%
  \setlength{\itemsep}{0.16em}%
  \setlength{\parsep}{0.1em}%
}{\end{enumerate}}
\newenvironment{examinstructions}{%
  \par\begingroup\itshape
}{%
  \par\endgroup\vspace{0.18em}
}
\makeatletter
\renewcommand\section{\@startsection{section}{1}{0pt}%
  {-1.25ex plus -.25ex minus -.1ex}%
  {0.55ex plus .1ex}%
  {\normalfont\large\bfseries}}
\renewcommand{\@maketitle}{%
  \newpage\null\vskip -1em%
  {\footnotesize\bfseries
    \href{https://www.ufl.edu/}{University of Florida}, \href{https://math.ufl.edu/}{Department of Mathematics}\par}%
  \vskip -0.5em%
  \tagstructbegin{tag=Title}%
    {\LARGE\bfseries\href{https://gma.math.ufl.edu/exams/topology-qualifying/}{Topology Qualifying Exam}, January 2026, Part 1\par}%
  \tagstructend%
  \vskip 0.25em%
  \tagmcbegin{artifact}%
    \hrule height 1pt%
  \tagmcend%
  \vskip 1em%
}
\makeatother
\title{Topology Qualifying Exam, January 2026, Part 1}
\author{}
\date{}
\begin{document}
\maketitle
\thispagestyle{empty}
\begin{examinstructions}
Be neat, do not write too small. Show your work and support your statements.
\end{examinstructions}
\begin{examproblems}{0}
\item
This problem has two parts.
\begin{examsubparts}
\item[\textnormal{(a)}]
Let \(f,g:X\to Y\) be continuous maps. Define what it means for~\(f\) to be homotopic to~\(g\).
\item[\textnormal{(b)}]
Suppose that the continuous maps \(f,g:X\to Y\) are homotopic and that the continuous maps \(h,k:Y\to Z\) are homotopic. Show that \(h\circ f\) is homotopic to \(k\circ g\).
\end{examsubparts}
\item
If \(j:A\hookrightarrow X\) is the inclusion of a retract, then show that the induced map on fundamental groups is injective.
\item
This problem has two parts.
\begin{examsubparts}
\item[\textnormal{(a)}]
State the Seifert–van Kampen theorem.
\item[\textnormal{(b)}]
Determine the fundamental group of \(S^1\vee S^1\).
\end{examsubparts}
\item
Prove that there are infinitely many distinct homotopy types of spaces with one 0-cell, one 1-cell, and one 2-cell.
\item
Let \(\mathbb{RP}^2\) denote the real projective plane.
\begin{examsubparts}
\item[\textnormal{(a)}]
Give the universal cover of \(\mathbb{RP}^2\) including the covering map.
\item[\textnormal{(b)}]
What is the group, \(G\), of covering transformations of this cover?
\item[\textnormal{(c)}]
How does it act on the universal cover?
\item[\textnormal{(d)}]
Does \(\mathbb{RP}^2\) have a cover, \(p:E\to\mathbb{RP}^2\), whose fundamental group, \(\pi_1(E)\), is \(\mathbb{Z}/(3)\)? Explain why or why not.
\item[\textnormal{(e)}]
Does \(\mathbb{RP}^2\) have a cover, \(p':E'\to\mathbb{RP}^2\), whose covering transformation group, \(G(E')\), is \(\mathbb{Z}/(5)\)? Explain why or why not.
\end{examsubparts}
\item
Give the definition of a chain complex of abelian groups and its homology groups.
\item
Let \((C,d)\), and \((C',d')\) be chain complexes of abelian groups.
\begin{examsubparts}
\item[\textnormal{(a)}]
Define what it means for \(f:(C,d)\to(C',d')\) to be a morphism of chain complexes, also called a chain map.
\item[\textnormal{(b)}]
For \(n\in\mathbb{Z}\) define the induced map \(H_n(f):H_n(C,d)\to H_n(C',d')\).
\item[\textnormal{(c)}]
Prove that \(H_n(f)\) is well defined.
\end{examsubparts}
\item
For \(n\geq 0\), let \(\Delta^n=\{(x_0,\ldots,x_n)\in\mathbb{R}^{n+1}\mid x_i\geq 0,\ \sum_i x_i=1\}\).
\begin{examsubparts}
\item[\textnormal{(a)}]
Let~\(X\) be a topological space. Define the singular chain complex \(\mathbf{S}_\bullet(X)\).
\item[\textnormal{(b)}]
Let~\(X\) be the one-point space \(*\). Compute \(H_n(X)\) for all~\(n\).
\end{examsubparts}
\item
Use the cellular chain complex to compute the homology of the Klein bottle.
\item
This problem has two parts.
\begin{examsubparts}
\item[\textnormal{(a)}]
State the Mayer–Vietoris sequence for reduced singular homology. You do not need to define the connecting homomorphism.
\item[\textnormal{(b)}]
Let \(n\geq 0\). Compute the reduced singular homology groups \(\widetilde{H}_k(S^n)\) for all~\(k\).
\end{examsubparts}
\end{examproblems}
\end{document}
